\section{Introduction}

Semantic-ID (\sid) tokenization has become a practical design point for
generative recommendation. Systems such as TIGER map items into discrete code
sequences so that a sequence model can generate candidate identifiers rather
than score every catalog item directly~\cite{rajput2023tiger}. This shift has
made the tokenizer itself a first-order artifact. Recent work explores
residual quantization frameworks, recommender-native encoding, differentiable
\sid learning, non-uniform quantization, qualified collisions, and continual
tokenization~\cite{ju2025grid,liang2026resid,fu2026diger,wei2026card,hu2026quasid,feng2026dact}.
The broader identifier literature also shows that item indexing choices matter
before generative retrieval is trained end-to-end~\cite{hua2023indexids}, that
semantic IDs can improve industrial ranking generalization~\cite{singh2023bettersemanticids},
and that collaborative tokenization methods explicitly optimize behavioral
alignment and assignment diversity~\cite{zhu2024cost,wang2024letter}.

The evaluation practice has not caught up with this artifact boundary. Prior
recommender tooling has argued for behavioral tests beyond NDCG-style
aggregates~\cite{chia2022reclist}, but SID tokenizers introduce a distinct
artifact: the generated item address space. A
tokenizer can support a usable downstream Recall@K while still hiding failure
modes that are hard to diagnose from aggregate ranking metrics: underused
codebooks, repeated full codes, semantically neat but behaviorally poor
prefixes, tail items sharing too little capacity, or prefix structures that
increase serving cost. These issues matter because \sid{}s are not merely
features. They define the address space exposed to the generator.

We present \tool, a resource for auditing item-to-\sid{} tokenizer artifacts
before or alongside downstream recommendation evaluation. The goal is
deliberately narrower than proposing a new tokenizer. \tool asks a concrete
question: given a public tokenizer artifact that maps items to code sequences,
what can be said about its capacity use, collisions, behavioral alignment,
head-tail allocation, and deployment cost without retraining a full generator?

The resource contribution is threefold. First, \tool defines a small
mapping-first interface for \sid assignments, item metadata, interaction
histories, and optional generator outputs. Second, it implements D1--D5a
diagnostics for utilization, collisions, semantic-collaborative alignment,
head-tail capacity, and structural serving-cost proxies, with D6 churn support
for continual-tokenizer settings. Third, it provides public resource-demo
evidence on GRID/RQ-KMeans-style, ReSID/GAOQ, and controlled sanity artifacts,
including a same-item Musical case study. The paper follows the CIKM Resource
four-page constraint, where appendices count toward the limit and detailed
logs are therefore placed in the online artifact~\cite{cikm2026resource}.
